\begin{definition}[Матрица смежности]
Для графа с фиксированным порядком вершин \(v_1,\ldots,v_n\) матрица
смежности \(A=(a_{ij})\) задаётся равенством \(a_{ij}=1\), если \(v_i\)
и \(v_j\) смежны, и \(a_{ij}=0\) иначе. У орграфа \(a_{ij}=1\)
означает наличие дуги \(v_i\to v_j\).
Для графа с кратными рёбрами \(a_{ij}\) равно их числу; для орграфа
с петлями \(a_{ii}\) равно числу петель. Далее неориентированные графы
с петлями в матричных формулах не рассматриваются.
\end{definition}
У неориентированного графа \(A=A^{\mathsf T}\); у простого графа
\(a_{ii}=0\), \(d(v_i)=\sum_j a_{ij}\). В орграфе суммы строк дают
исходящие, а суммы столбцов — входящие степени.
\begin{graphillustration}{Пример: граф $P_3$ и его матрица в порядке $v_1,v_2,v_3$.}{fig:01-adjacency}
\begin{tikzpicture}[font=\small]
\node[gi vertex] (v1) at (0,0) {$v_1$};
\node[gi vertex] (v2) at (1.5,1) {$v_2$};
\node[gi vertex] (v3) at (3,0) {$v_3$};
\draw[gi edge] (v1) -- (v2);
\draw[gi edge] (v2) -- (v3);
\node[gi note] at (6,0) {$A=\begin{pmatrix}0&1&0\\1&0&1\\0&1&0\end{pmatrix}$};
\node[gi note] at (6,-1.25) {Строка $i$ — соседи $v_i$};
\end{tikzpicture}
\end{graphillustration}


\begin{lemma}[Разбиение по предпоследней вершине]
Пусть \(w_k(i,j)\) — число маршрутов длины \(k\) из \(v_i\) в \(v_j\).
Тогда
\[
 w_0(i,j)=\begin{cases}1,&i=j,\\0,&i\ne j,\end{cases}
 \qquad w_{k+1}(i,j)=\sum_{s=1}^n w_k(i,s)a_{sj}.
\]
\end{lemma}
\begin{proof}
Маршрут нулевой длины состоит из одной вершины. У каждого маршрута
длины \(k+1\) однозначно определены первые \(k\) рёбер и последнее ребро.
Если предпоследняя вершина \(v_s\), выборов этих частей соответственно
\(w_k(i,s)\) и \(a_{sj}\). Разные значения \(s\) дают непересекающиеся классы.
\end{proof}
\begin{theorem}[Комбинаторный смысл степеней матрицы]
Для любого целого \(k\ge0\) элемент \((A^k)_{ij}\) равен числу маршрутов
длины \(k\) из \(v_i\) в \(v_j\).
\end{theorem}
\begin{proof}
Индукция по \(k\). При \(k=0\) имеем \(A^0=I\), что совпадает
с \(w_0(i,j)\). Если утверждение верно для \(k\), то
\[
 (A^{k+1})_{ij}=\sum_s(A^k)_{is}a_{sj}
 =\sum_s w_k(i,s)a_{sj}=w_{k+1}(i,j)
\]
по лемме. Индукционный переход доказан.
\end{proof}
\begin{cor}
Вершина \(v_j\) достижима из \(v_i\) тогда и только тогда, когда
\((I+A+\cdots+A^{n-1})_{ij}>0\).
\end{cor}
\begin{proof}
Положительное слагаемое даёт маршрут, значит, достижимость.
Обратно, из маршрута удалением повторений получается простой путь;
при различных концах в нём не более \(n-1\) рёбер. При совпадающих
концах достаточно маршрута длины ноль.
\end{proof}
\begin{example}[Почему это не число простых путей]
В графе из одного ребра \(uv\) существует маршрут \(u,v,u\) длины два,
поэтому \((A^2)_{uu}=1\). Простого пути длины два из \(u\) в \(u\) нет.
В лекциях такой маршрут называется путём; простота пути оговаривается отдельно.
\end{example}
